% WakeKV workshop-paper draft — Phase 0-2a snapshot.
% RAW DRAFT — kept in-repo, not for release. Real numbers filled in;
% prose that needs a real pass is marked \NEEDPROSE{...}. Figures/tables
% referenced but not yet generated are marked \NEEDFIG.
%
% Compile: pdflatex wakekv.tex  (twice for refs). Uses only base + amsmath +
% booktabs + hyperref + graphicx + a lightweight geometry setup so it
% builds without a specific workshop's style file yet.

\documentclass[11pt]{article}
\usepackage[letterpaper,margin=1in]{geometry}
\usepackage{amsmath,amssymb}
\usepackage{booktabs}
\usepackage{graphicx}
\usepackage{hyperref}
\usepackage{xcolor}
\usepackage{url}

\newcommand{\NEEDPROSE}[1]{\textcolor{red}{[TODO: #1]}}
\newcommand{\NEEDFIG}{\textcolor{red}{[FIG PENDING]}}

\title{When Do Attention Heads Change Their Minds?\\
       A Case for Reactive KV-Cache Residency}
\author{Utkarsh Ranjan\\ UC San Diego}
\date{Draft --- \today}

\begin{document}
\maketitle

\begin{abstract}
Recent KV-cache optimizations classify attention heads once --- offline
(DuoAttention, FlexiCache) or at prefill (Ada-KV, LAVa, HeteroCache) ---
and then treat the classification as fixed for the rest of the
generation. We measure per-decode-step head activity across three
regimes (needle retrieval, long chain-of-thought, multi-turn recall) on
Qwen2.5-3B, DeepSeek-R1-Distill-Qwen-1.5B, and DeepSeek-R1-Distill-Llama-8B
and find 63--85\% of heads shift their reading habits at least once
within a single generation --- including at 8B scale --- replicating
the retrieval-heads-are-dynamic phenomenon on our hardware.
We then ask whether cheap runtime signals can \emph{predict} wake-ups
with enough lead time to prefetch the head's KV cache from CPU. Under a
fair evaluation (one global threshold per signal, causal per-head
normalization) they cannot: precision $0.06$--$0.20$ at recall
$0.41$--$0.72$ across four candidates. Wake-ups are temporally bursty
(mean concentration $z = +9.7$ vs. shuffle null on CoT) but not tightly
detectable enough for a coarse trigger either. We therefore argue for a
\emph{reactive} design --- demote cooling heads to CPU and fetch on
demand --- and validate it, at page granularity over the same logged
attention, against two competitors: a FlexiCache-style frozen policy and
a ReasonAlloc-style policy that shares reactive's exact LRU schedule but
demotes destructively rather than reversibly. Across five real
model/regime combinations (three model families, 1.5B--7B, three
regimes), reactive beats frozen at every comparable matched-memory
operating point (17/17) and beats destructive demotion at matched budget
in 19 of 20 settings --- isolating reversibility itself, not just
dynamism, as a source of the advantage. Preliminary real-system
measurements on the A30 (Mistral-7B via a FlexiCache/vLLM shim) confirm
the frozen-vs-reactive direction: reactive residency achieves
1.34--2.33$\times$ throughput at 95--99\% LongBench mean quality vs.\
stock FlexiCache across the rerank-interval sweep. Broader real-system
evaluation (larger models, multi-turn, a real-system reversible-vs-evict
head-to-head) is future work; the negative prediction result and the
simulator-level Pareto and reversibility results stand on their own.
\end{abstract}

\section{Introduction}
\label{sec:intro}

\NEEDPROSE{one-paragraph hook: KV-cache size dominates long-context
serving memory; the dominant response is to compress or evict per-head,
under one uniform assumption --- that each head's role during a
generation can be summarized once. Recent measurement work on retrieval
heads (\cite{retrievalheadsdynamic}) casts doubt on that assumption at
the retrieval-behavior level. This paper asks what that means for the
cache.}

Our contributions are:
\begin{enumerate}
    \item A replication of head non-stationarity on our own models and
        regimes, with two novel supporting regimes (long CoT, scripted
        multi-turn recall) that competing frozen-role methods do not
        evaluate (\S\ref{sec:churn}).
    \item An \emph{honest negative} result: cheap per-head signals, and
        an ensemble of them, cannot predict wake-ups well enough to
        prefetch under a fair evaluation; a structurally different
        signal (generated-token identity, discovered without a
        hand-picked vocabulary) does clear a useful bar, but only in
        long-CoT reasoning traces (\S\ref{sec:predict}).
    \item A characterization of wake-up burstiness --- real but not
        tightly detectable --- ruling out coarse boundary triggers as
        well (\S\ref{sec:cluster}).
    \item A page-granularity residency simulator showing reactive LRU +
        fetch-on-demand dominates FlexiCache-style frozen classification
        on the memory/miss Pareto (C2), and that reversible offload
        specifically --- not just LRU dynamism --- beats destructive
        eviction at matched budget (C3), across five real model/regime
        combinations (\S\ref{sec:sim}).
\end{enumerate}

\textbf{Explicit scope.} The paper's primary contributions are the
measurement of head non-stationarity, the negative prediction result,
and the simulator-level reactive-vs-frozen and reversible-vs-evict
results (contributions 1--4 above). A preliminary real-system evaluation
(\S\ref{sec:limits},
Table~\ref{tab:m2b2}) confirms the direction on physical hardware but
is narrow in scope: single model, single input/output regime, small
sample count, one A30. A full serving-system evaluation across
model scales, workloads, and hardware is future work.

\section{Setup}
\label{sec:setup}

\textbf{Models.} Qwen2.5-3B-Instruct (36 layers $\times$ 16 heads = 576
attention heads), DeepSeek-R1-Distill-Qwen-1.5B (28 layers $\times$ 12
heads = 336), and DeepSeek-R1-Distill-Llama-8B (32 layers $\times$ 32
heads = 1024, GQA with 8 KV-groups, bf16 on Ampere). Two hardware
epochs: measurements on 1.5B/3B were taken on Turing (RTX~2080~Ti, 11~GB, fp16 or float32
where fp16 overflows); the 8B measurements were taken on Ampere
(A30, 24~GB, bf16, PCIe~Gen4).

\textbf{Regimes.} (i) NIAH, a synthetic UUID needle in $\approx 5000$
tokens of filler; (ii) long CoT, math problems from MATH-500 with up to
2048 generated tokens on the reasoning distillation; (iii) multi-turn,
a scripted conversation with facts stated in turn 1, three unrelated
turns of filler, and a final recall request.

\textbf{Instrumentation.} We modify the transformers decode loop to log,
at every decode step and every attention head, the top-$k$ attended
positions and their weights ($k = 64$), a continuous needle score
following \cite{retrievalheadsdynamic}'s Eq.~2, and a binary copy-paste
score (Eq.~1). To fit within 11~GB of GPU memory we never materialize
the full attention map: prefill runs without attention output, and each
decode step's per-head attention row is reduced to top-$k$ immediately.
One run produces $\approx 100$--$200$~MB of logs.

\textbf{Hardware.} The 1.5B/3B measurements were taken on a single
NVIDIA RTX~2080~Ti (11~GB, Turing) --- scout-scale by design. The
8B measurements were taken on a single NVIDIA A30 (24~GB, Ampere,
PCIe Gen4 $\times$16, $\approx$25~GB/s achieved H2D). The follow-up
serving-system evaluation (\S\ref{sec:limits}) uses the A30 throughout.

\section{Do Heads Shift During Decoding?}
\label{sec:churn}

We ask whether the head set responsible for retrieving from context
changes step-to-step within a single generation, and characterize the
change with three complementary metrics.

\textbf{Adjacent-step Jaccard.} Following \cite{retrievalheadsdynamic},
for each step $t$ we identify the set of heads whose copy-paste score
fires at $t$ and compute the Jaccard similarity to the set at $t{+}1$.
Under the continuous needle score at threshold 0.1, the mean over NIAH
runs is $0.53$ ($\pm 0.06$), squarely within the $0.28$--$0.51$ range
reported by \cite{retrievalheadsdynamic}, replicating the phenomenon on
our hardware.

\textbf{FlexiCache-style RCO temporal stability.} For each head, we
compute the random-corrected top-$k$-page overlap between decode step
$t$ and $t + \Delta$ for $\Delta \in [1, 15]$, and average. The fraction
of heads with mean $\mathrm{RCO} < 0.5$ (``chronically unstable'') is
$0.02$--$0.09$ across regimes, consistent with HeteroCache's $\sim$7\%
``volatile'' class, independently reproduced.

\textbf{Drift events.} Following HeteroCache's trigger, we detect
sustained overlap drops against a rolling baseline (windowed median
$< 0.5$). The fraction of heads with at least one drift event within a
single generation ranges from $0.63$ (NIAH, 64-step run) to
$0.85$ (long CoT).

The three metrics agree on structure and disagree on magnitude in the
expected way: a small tail of chronically unstable heads coexists with
a large majority of occasionally-shifting heads. Table~\ref{tab:churn}
summarizes.

\begin{table}[t]
\centering
\small
\begin{tabular}{lrrrr}
\toprule
Regime & Model & Adj.~Jaccard & Unstable frac.\ & $\geq$1 drift event \\
\midrule
NIAH & Qwen2.5-3B & 0.53 (cont.) & 0.064 & 0.63 \\
CoT & R1-Distill-1.5B & 0.82 (bin.) & 0.023 & 0.85 \\
Multi-turn & Qwen2.5-3B & 0.79 (bin.) & 0.085 & 0.76 \\
\midrule
CoT & \textbf{R1-Distill-8B} & \textbf{0.93 (bin.)} & \textbf{0.042} & \textbf{0.79} \\
\bottomrule
\end{tabular}
\caption{Head-shift metrics by regime. Adjacent Jaccard reported under
the continuous score (NIAH) and binary score (CoT / multi-turn); see
\S\ref{sec:limits} for the artifact caveat. Unstable fraction uses
FlexiCache's temporal-stability metric with threshold 0.5; drift
fraction uses HeteroCache's trigger. At 8B (bottom row), the binary
Jaccard tightens to 0.93 as the model's copy behavior stabilizes on
math CoT, but the page-level drift metric --- the signal relevant to
a residency controller --- remains high at 0.79, so the residency
case is undiminished at scale.}
\label{tab:churn}
\end{table}

\NEEDFIG~Figure~\ref{fig:heatmap}: per-step per-head activity heatmap
on a representative CoT run, with the twelve most-active heads
highlighted (analogous to the ``retrieval heads are dynamic'' Fig.~1).

\section{Can We Predict the Shifts Cheaply?}
\label{sec:predict}

If we could \emph{predict} that a currently-quiet head is about to become
active a few decode steps ahead, we could prefetch its KV pages from CPU
and hide the promotion latency. We test four candidate signals:
\texttt{online\_rco} (change in the head's top-$k$ page set),
\texttt{drift} (HeteroCache's windowed-median overlap), \texttt{entropy\_trend}
(rise in attention entropy), and \texttt{needle\_mass\_delta}
(an oracle-ish reference).

\textbf{Wake events.} A head is said to ``wake'' when its continuous
needle score crosses from below a low threshold to above a high threshold
with a minimum quiet period (hysteresis, to avoid classifying jitter as
wake-ups). This is the ground truth we grade signals against.

\textbf{The honesty ladder.} We grade each signal three ways:

\begin{enumerate}
    \item \emph{Per-head best threshold} (oracle) --- picks a cutoff per
        head after seeing labels. Under this grading, \texttt{drift}
        appears perfect (P=1.00, R=1.00 at lead 8 on CoT). This is not
        deployable.
    \item \emph{Single fixed global threshold} --- one raw cutoff for all
        heads. Precision-starved: no signal clears a useful bar. Best
        F1: $0.15$ (CoT).
    \item \emph{Causal per-head z-score} --- each head normalized against
        its own trailing 32-step window, then one global $z > 2$
        threshold. This is what a runtime controller can compute; it is
        the honest number.
\end{enumerate}

Table~\ref{tab:predict} reports the deployable numbers. The best signal
(\texttt{drift}) achieves precision $0.08$--$0.20$ at recall
$0.41$--$0.70$: over $80\%$ of alarms are false while $30$--$60\%$ of
real wake-ups are missed. A precision this low cannot be tuned to a
usable operating point --- and for WakeKV specifically it bites twice:
missed wake-ups become promotion stalls, and false alarms
over-promote heads to the GPU, evaporating the memory savings the
system was built for.

\begin{table}[t]
\centering
\small
\begin{tabular}{lrrrr}
\toprule
Signal & Regime & Precision & Recall & F1 \\
\midrule
drift        & CoT 1.5B (407 events) & 0.08 & 0.41 & 0.14 \\
entropy\_trend & CoT 1.5B             & 0.06 & 0.45 & 0.11 \\
online\_rco  & CoT 1.5B               & 0.06 & 0.69 & 0.11 \\
\midrule
drift        & Multi-turn 3B (73)    & 0.20 & 0.70 & 0.31 \\
entropy\_trend & Multi-turn 3B       & 0.20 & 0.56 & 0.29 \\
online\_rco  & Multi-turn 3B         & 0.13 & 0.67 & 0.22 \\
\midrule
\textbf{online\_rco}   & \textbf{CoT 8B} & \textbf{0.10} & \textbf{0.79} & \textbf{0.17} \\
drift                  & CoT 8B          & 0.05 & 0.44 & 0.10 \\
needle\_mass\_delta    & CoT 8B          & 0.04 & 0.79 & 0.10 \\
entropy\_trend         & CoT 8B          & 0.04 & 0.60 & 0.08 \\
\bottomrule
\end{tabular}
\caption{Deployable P/R/F1 at lead 32 under the causal per-head z-score
grading ($z > 2$, trailing window 32). Every candidate signal is
precision-starved at every scale. The best-signal ranking is
scale-dependent (\texttt{drift} wins at 1.5B/3B, \texttt{online\_rco}
wins at 8B), but none crosses a usable operating point.}
\label{tab:predict}
\end{table}

\NEEDFIG~Figure~\ref{fig:predict}: precision-recall curves at lead 32
for the four signals on CoT, showing all four hug the low-precision
corner.

\textbf{Two follow-up candidates.} Before committing to the reactive
design, we tested two structurally different ideas the four
attention-derived signals above don't cover.

\emph{Ensemble voting} combines the four causally z-scored signals,
firing when $\geq 2$ of 4 agree. This is a negative result: across all
four regimes tested (multi-turn, NIAH, CoT at 1.5B and 8B), the ensemble
never beats the single best individual signal (F1 $0.17$ vs.\ drift's
$0.31$ on multi-turn; $0.18$ vs.\ $0.26$ on NIAH; $0.09$ vs.\ $0.10$ on
CoT~8B; $0.13$ vs.\ $0.14$ on CoT~1.5B). The four signals' false
positives are not decorrelated enough for agreement to help --- three of
the four are all derived from top-$k$ page-overlap statistics, so they
tend to fire together or not at all.

\emph{Token/wake correlation} asks whether the identity of the just-generated
token --- not any head's attention --- predicts a wake burst. Rather than
positing a hand-picked marker vocabulary (an uncited, arbitrary guess that
would be circular if tuned against results), we scan every generated
token id, test each with a two-proportion $z$-test against the run's base
wake rate, Bonferroni-correct across every token tested, and evaluate the
surviving markers on a held-out half of the data disjoint from what
selected them. The result is regime-specific: zero tokens survive the
Bonferroni bar on multi-turn or NIAH, but on both CoT runs (DeepSeek-R1-Distill,
reasoning traces) the discovered markers reach held-out F1 $0.53$--$0.62$
at lead 32 --- $4$--$5\times$ every other signal in this paper on the
same data, and, because of the discovery/held-out split, an honest number
rather than an oracle-tuned one. This supports rather than contradicts
the reactive design: it is CoT-specific (token identity is a
model-vocabulary-specific signal needing per-model recalibration, unlike
the architecture-generic attention statistics above), recall is still
$0.68$--$0.89$ (not $1.0$, so reactive fallback stays mandatory
regardless), and precision $0.43$--$0.48$ still means roughly half of
triggered prefetches would be unnecessary. We report it as evidence that
we searched the design space thoroughly --- including the most
promising direction, generated-token identity --- before committing to
reactive, not as a mechanism WakeKV deploys.

\section{Are Wake-ups Clustered in Time?}
\label{sec:cluster}

If wake-ups happen in bursts --- many heads waking together at reasoning
phase boundaries, for instance --- a single coarse trigger could
prefetch cheaply where the per-head prediction of \S\ref{sec:predict}
failed. We test this by pooling per-head wake events into a per-step
histogram, then comparing the concentration in the busiest 10\% of
steps to a per-head uniform-shuffle null (200 trials, preserving each
head's event count).

On the CoT regime, mean concentration is $z = +9.7$ against the null,
significant in 4 of 5 runs (one run reaches $z = +29.4$ with 22 heads
waking at a single step). Wake-ups are genuinely bursty --- \emph{not}
Poisson.

But the clustering is not tight enough to give a coarse trigger a
clean operating point. A trigger firing at the busiest 2\% of steps
($\pm 4$ step window) catches only 54\% of wake-ups. Widening to 5\%
of steps catches 86\% but the concentration advantage over the null
shrinks to $1.29\times$ (5\% of a long reasoning trace is a big chunk).
And in either configuration, the ``busiest steps'' are chosen after
seeing the events: at runtime you would need a detectable boundary
signal (a transition token, a global entropy spike) that aligns with
the bursts, which we have not shown exists.

Both cheap proactive routes --- per-head prediction and coarse
boundary triggers --- are therefore ruled out.

\NEEDFIG~Figure~\ref{fig:cluster}: per-step wake-up histogram for
math500-1 (the highest-$z$ CoT run) with the shuffle-null envelope
overlaid.

\section{Reactive Residency Wins on the Memory/Miss Pareto}
\label{sec:sim}

Given prediction fails, we validate a reactive design: demote cooling
pages to a CPU reservoir, fetch on demand when a head accesses a
demoted page. We simulate four policies at page granularity over the
same logged attention:

\begin{description}
    \item[Full] every seen page stays GPU-resident (quality and memory
        ceiling, zero misses).
    \item[Frozen] FlexiCache-style: heads classified stable/unstable
        once, offline, from their own churn. Unstable heads keep full
        residency; stable heads keep a fixed top-$B$ set refreshed
        every 16 steps. Classification never changes during decode.
    \item[Reactive (WakeKV)] every head capped at $B$ resident pages
        via LRU; a wanted page in the CPU reservoir is fetched on
        demand and promoted, evicting the least-recently-wanted page
        back to the reservoir. Nothing is destroyed.
    \item[Evict (ReasonAlloc-style)] identical LRU cap and eviction
        \emph{schedule} to Reactive, same budget $B$ --- but demotion is
        destructive: an evicted page has no reservoir, so a later want
        of it misses again, forever. This isolates \emph{reversibility}
        itself as a variable, holding dynamism, budget, and eviction
        order fixed on both sides.
\end{description}

\textbf{C2 result: reactive beats frozen.} Table~\ref{tab:sim}
interpolates reactive's miss rate to each frozen operating point's mean
memory (a matched-memory comparison, since at the same budget label the
two policies sit at different memory). Reactive dominates frozen at
every comparable operating point across all five model/regime combos
tested (17/17), at roughly half the miss rate on the smaller models and
up to $\sim$$10\times$ on Mistral-7B/NIAH.

\begin{table}[t]
\centering
\small
\begin{tabular}{llrrr}
\toprule
Regime & Model & Mean mem (pages) & Frozen miss & Reactive miss \\
\midrule
CoT       & 1.5B & 4980 & 0.395 & 0.216 \\
CoT       & 1.5B & 6627 & 0.191 & 0.121 \\
CoT       & 1.5B & 8459 & 0.066 & 0.034 \\
CoT       & 1.5B & 9233 & 0.042 & 0.024 \\
\midrule
Multi-turn & 3B & 7918  & 0.491 & 0.392 \\
Multi-turn & 3B & 11039 & 0.294 & 0.234 \\
Multi-turn & 3B & 15828 & 0.116 & 0.065 \\
Multi-turn & 3B & 16042 & 0.110 & 0.057 \\
\midrule
NIAH & 3B & 34047 & 0.403 & 0.086 \\
\midrule
\textbf{CoT} & \textbf{8B} & \textbf{18388} & \textbf{0.455} & \textbf{0.269} \\
\textbf{CoT} & \textbf{8B} & \textbf{23749} & \textbf{0.257} & \textbf{0.149} \\
\textbf{CoT} & \textbf{8B} & \textbf{31372} & \textbf{0.102} & \textbf{0.060} \\
\textbf{CoT} & \textbf{8B} & \textbf{37321} & \textbf{0.045} & \textbf{0.025} \\
\midrule
NIAH & Mistral-7B & 35191 & 0.578 & 0.164 \\
NIAH & Mistral-7B & 38763 & 0.504 & 0.146 \\
NIAH & Mistral-7B & 45713 & 0.441 & 0.110 \\
NIAH & Mistral-7B & 56297 & 0.405 & 0.055 \\
\bottomrule
\end{tabular}
\caption{Matched-memory miss rates from the residency simulator, all
five real model/regime combos tested. Reactive is interpolated onto
each frozen operating point's mean resident memory; NIAH/3B has only
one comparable point because reactive's memory footprint stayed below
all but the smallest frozen budget in that run, not a partial result.
Reactive wins 17/17 comparable points. Mistral-7B is the same model
the preliminary real-system evaluation (\S\ref{sec:limits}) uses ---
the first model with both simulator and real-system coverage, though
its own churn (\S\ref{sec:churn}) has not separately been measured.}
\label{tab:sim}
\end{table}

\textbf{Miss asymmetry --- the honest framing.} A reactive miss and a
frozen miss are not the same event. A reactive miss is a \emph{paid
fetch stall}: the page is loaded from CPU, decode blocks briefly, and
quality is preserved. A frozen miss is a \emph{quality gap}: the page
is unavailable until the next scheduled refresh. So reactive trades
memory for stall time, not for accuracy. Whether the trade is
favorable end-to-end depends on the stall time relative to the memory
saved, which the simulator cannot measure; this is the central
question the follow-up (\S\ref{sec:limits}) addresses.

\NEEDFIG~Figure~\ref{fig:pareto}: memory-vs-miss Pareto curves for
frozen and reactive across the sweep of $B$, for both regimes.

\textbf{C3 result: reversibility itself, isolated.} Reactive and Evict
share the identical LRU rule and budget, so no memory-matching
interpolation is needed --- any miss-rate gap is attributable to
demotion consequence alone, not to a confound in dynamism or budget.
Table~\ref{tab:evict} reports reactive vs.\ evict at the same budget
across all five combos. Reactive wins 19/20 matched budgets; the one
exception is CoT/8B at the tightest budget ($B{=}8$: reactive $0.632$
vs.\ evict $0.624$, evict slightly better). We attribute this to a
real second-order effect rather than noise in the comparison: the two
policies' eviction \emph{schedules} are only provably identical up to
the first page reactive recovers that evict left destroyed --- after
that, reactive's resident set literally contains content evict's
doesn't, so a later eviction round can legitimately pick a different
page under each policy. At the tightest budget on the highest-churn
regime, recovery events are most frequent, so this composition drift
has the most room to compound, which is exactly the one place it
surfaces. It does not undermine the 19/20 result.

\begin{table}[t]
\centering
\small
\begin{tabular}{llrrr}
\toprule
Regime & Model & Budget $B$ & Reactive miss & Evict miss \\
\midrule
Multi-turn & 3B & 8  & 0.631 & 0.651 \\
Multi-turn & 3B & 16 & 0.300 & 0.369 \\
Multi-turn & 3B & 32 & 0.010 & 0.025 \\
Multi-turn & 3B & 64 & 0.000 & 0.000 \\
\midrule
NIAH & 3B & 8  & 0.655 & 0.760 \\
NIAH & 3B & 16 & 0.409 & 0.617 \\
NIAH & 3B & 32 & 0.187 & 0.398 \\
NIAH & 3B & 64 & 0.076 & 0.215 \\
\midrule
\textbf{CoT} & \textbf{1.5B} & \textbf{8}  & \textbf{0.542} & \textbf{0.543} \\
\textbf{CoT} & \textbf{1.5B} & \textbf{16} & \textbf{0.221} & \textbf{0.263} \\
\textbf{CoT} & \textbf{1.5B} & \textbf{32} & \textbf{0.040} & \textbf{0.074} \\
\textbf{CoT} & \textbf{1.5B} & \textbf{64} & \textbf{0.005} & \textbf{0.013} \\
\midrule
CoT & 8B & 8  & 0.632 & \textbf{0.624 (evict better)} \\
CoT & 8B & 16 & 0.338 & 0.373 \\
CoT & 8B & 32 & 0.090 & 0.145 \\
CoT & 8B & 64 & 0.014 & 0.035 \\
\midrule
NIAH & Mistral-7B & 8  & 0.650 & 0.761 \\
NIAH & Mistral-7B & 16 & 0.443 & 0.613 \\
NIAH & Mistral-7B & 32 & 0.180 & 0.309 \\
NIAH & Mistral-7B & 64 & 0.048 & 0.091 \\
\bottomrule
\end{tabular}
\caption{Reactive (offload) vs.\ evict (destroy) at the same budget,
all five real model/regime combos. No interpolation --- both policies
share the identical LRU cap and eviction order, so demotion consequence
is the only variable. 19/20 matched budgets favor reversibility; the
sole exception (bold row) is explained in the main text. This is the
first real-log confirmation of C3; previously it had only a
synthetic-data check.}
\label{tab:evict}
\end{table}

\section{Related Work}
\label{sec:related}

\NEEDPROSE{Expand the compact paragraph below into two or three
paragraphs organized by our positioning axes: (1) static head-role
methods, (2) prefill-only adaptive allocators, (3) decode-time
compressors, (4) offload systems. WakeKV closes the intersection
between (4)'s reversibility and dynamic per-head adjustment.}

Static head-role methods (DuoAttention, RazorAttention, MoA, HeadKV,
FlexiCache, HeteroCache) fix each head's cache treatment offline or at
prefill. Adaptive budget allocators (Ada-KV, CAKE, LAVa, DynamicKV) are
prefill-only despite their names. Decode-time compressors (R-KV, G-KV,
ThinKV, SCOPE) evict under uniform per-head budgets. ReasonAlloc
reallocates head budgets during decoding but demotion permanently
destroys KV. Offload systems (ArkVale, ShadowKV, InfiniGen, HeadInfer,
FlexiCache, HeteroCache) keep KV recoverable but never change head
priority at runtime. \cite{retrievalheadsdynamic} shows head roles
churn within a single generation and names dynamic per-head KV
retention as an open direction. WakeKV closes the intersection.

\section{Limitations and Future Work}
\label{sec:limits}

We keep this section deliberately blunt because it is what a reviewer
should read before the introduction.

\textbf{Scale.} 1.5B and 3B measurements are on Turing (11~GB);
8B measurements are on Ampere (A30, 24~GB, PCIe~Gen4). The 8B validation
confirms the churn phenomenon (Table~\ref{tab:churn}), the deployable
prediction failure (Table~\ref{tab:predict}), and the reactive-vs-frozen
Pareto (Table~\ref{tab:sim}) all hold. However, the 8B run also shows
that reactive traffic scales harshly: the simulator counts millions of
fetches at low budgets (up to $\sim$19M at $B{=}8$), whose real cost
--- unmeasurable here --- Phase 2b must quantify. Nothing above 8B was
tested; a 14B/70B extension is future work.

\textbf{Sample size.} Five CoT runs (407 wake events, enough for the
Table~\ref{tab:predict} numbers to have real weight) but only one
multi-turn run (73 events, coarser).

\textbf{Simulator scope.} The residency simulator counts misses, not
stall time. It works at page granularity approximated from the logged
top-$k$, not the full KV cache. It measures no quality metric --- it
assumes fetch-on-demand preserves quality by construction, but only
running LongBench/RULER/SCBench confirms this. And it compares
against a FlexiCache-\emph{style} policy, not FlexiCache's actual
code.

\textbf{Binary-metric artifact.} The 2602.11162 binary retrieval score
is thresholded, so adjacent-step Jaccard reads higher under it than
under the continuous score. We report both.

\textbf{Preliminary real-system evaluation.} We measure the shim
described above on the A30 against a stock-FlexiCache baseline
(Mistral-7B-Instruct-v0.2, LEval-style 8k input / 1000 output, 24
prompts) on two axes: throughput and LongBench quality
(Table~\ref{tab:m2b2}).

\textbf{Throughput.} Reactive residency beats stock at every rerank
interval, monotone --- 1.34$\times$ at $R{=}1$ (the WakeKV design
point) up to 2.33$\times$ at $R{=}16$. This contradicted our initial
prediction that reactive's every-step rerank cost would trail stock at
low intervals, and has a mechanistic explanation: the shim removes
FlexiCache's offline classification entirely, so every head runs the
fast sparse-decode path (0 dense heads vs.\ 64 in stock). Part of the
win is doing less attention work per step, not purely better memory
management. Holding rerank interval fixed at the native cadence
($R{=}16$) isolates this directly, reusing rows already in
Table~\ref{tab:m2b2}: dropping from 64 dense heads to 0 more than
doubles throughput (2.33$\times$) for a 5-point LongBench cost ---
confirming sparsity alone already accounts for most of the throughput
headline and essentially all of the quality cost
(\texttt{notes/wakekv\_m2b2\_sweep\_a30.md}). This does not yet isolate
rerank frequency's own contribution, though: the complementary cell
(64 dense heads reranked every step) is still future work (below).

\textbf{Quality.} LongBench mean (4 tasks, 30 samples each) stays
95--99\% of stock across the whole rerank-interval sweep, with no
monotonic erosion --- even the fastest point (R=16, 2.33$\times$
throughput) holds 95\%. TriviaQA F1 is identical across all six
configurations; the spread is dominated by the two small-N QA tasks
(qasper, 2wikimqa) bouncing $\pm 2$--$4$ points, indicating per-task
deltas are noise at N=30. Reactive residency therefore buys throughput
without measurably trading quality on this benchmark at this scale.

\begin{table}[t]
\centering
\small
\begin{tabular}{lrrrr}
\toprule
Config & Rerank int.\ & Gen tok/s & vs stock & LB \% \\
\midrule
Stock FlexiCache & 16 (native) & 114.4 & 1.00$\times$ & 100 \\
Reactive (WakeKV) & 1 & 153.6 & 1.34$\times$ & 97.4 \\
Reactive & 2 & 193.3 & 1.69$\times$ & 97.1 \\
Reactive & 4 & 226.3 & 1.98$\times$ & 99.1 \\
Reactive & 8 & 254.2 & 2.22$\times$ & 97.9 \\
Reactive & 16 & 266.6 & 2.33$\times$ & 95.0 \\
\bottomrule
\end{tabular}
\caption{Preliminary real-system throughput$\times$quality sweep on
the A30 (Mistral-7B, LEval 8k$\to$1000, 24 prompts; LongBench 4
tasks $\times$ 30 samples). Reactive residency beats stock
FlexiCache at every rerank interval on throughput; LongBench mean
stays 95--99\% of stock across the whole sweep with no monotonic
erosion. The R=1 row is the WakeKV design point; R=16 is the
Pareto knee here (largest throughput win, smallest quality gap
still $\geq 95\%$). Data:
\texttt{notes/wakekv\_m2b2\_sweep\_a30.md}.}
\label{tab:m2b2}
\end{table}

\textbf{Scope of the preliminary result.} Single model
(Mistral-7B), single input/output regime (8k$\to$1000, 24 prompts),
30 LongBench samples per task, A30 with 24~GB and 16~GB host KV
pool, native FlexiCache LEval configuration. Ranks the two policies
on this specific configuration; does not claim serving performance
on other hardware or model scales. The LongBench mean is more
stable than the per-task numbers at this sample count.

\textbf{Future work.} A follow-up on the same stack will (i) extend
the quality evaluation to the full LongBench task list at
$\sim$200 samples/task, and to RULER and SCBench; (ii) cover the 8B
and multi-turn regimes where the Phase 2a simulator showed the
strongest memory/miss wins; (iii) measure real PCIe stall time from
the fetch-on-demand path (currently visible only through end-to-end
throughput); (iv) run the complementary sparse$\times$rerank cell ---
FlexiCache's 64-dense-head classification reranked every step,
matching WakeKV's native cadence --- to isolate rerank frequency's own
contribution from the dense-head-count effect already isolated above;
(v) run a real-system evict-vs-offload head-to-head (C3). The
simulator-level reading (\S\ref{sec:sim}, Table~\ref{tab:evict}) now
supports reversible offload on real logs across all five model/regime
combos (19/20 matched budgets). The controller-level code for a real
evict mode is written and tested (\texttt{ReactiveController}'s
\texttt{demotion="evict"}, validated against the simulator's numbers),
but wiring it into the FlexiCache shim specifically is not a simple
config override the way reactive mode is: tracing FlexiCache's actual
promote/demote path shows reversibility is an emergent property of
three cooperating structures (a generic GPU block free-list with no
memory of prior contents, an unconditional per-token CPU mirror, and a
separately-computed MinMax score cache), and the one clean patch point
found needs a new per-request state tensor threaded through four
request-lifecycle methods to avoid a completed request's eviction
history leaking into a reused batch slot --- not verifiable without a
live GPU run, so it remains a documented design (not yet implemented)
rather than a real-system result.

\section*{Reproducibility}
Source, logs, and results at \NEEDPROSE{GitHub URL after de-anonymize}.
All Phase 0--2a analyses can be re-run offline on the logs without a
GPU.

\begin{thebibliography}{9}
\bibitem{retrievalheadsdynamic}
    \NEEDPROSE{Retrieval Heads are Dynamic (arXiv:2602.11162), full citation.}
\NEEDPROSE{FlexiCache (arXiv:2511.00868, MLSys 2026)}
\NEEDPROSE{HeteroCache (arXiv:2601.13684, ACL 2026)}
\NEEDPROSE{ReasonAlloc (arXiv:2606.11164)}
\NEEDPROSE{Ada-KV, LAVa, DynamicKV, CAKE, HeadKV, DuoAttention, RazorAttention, MoA, ThinKV, R-KV, ArkVale, ShadowKV, InfiniGen, HeadInfer, SCBench, LongBench, RULER, HotpotQA/MATH-500 --- to be filled from notes/big\_four\_deep\_read.md.}
\end{thebibliography}

\end{document}
